If some submodule had no finite decomposition into irreducible submodules, the ascending chain condition would give a maximal such submodule $N$. It is not irreducible, so $N=N_1\cap N_2$ with both $N_i$ strictly larger. By maximality, each $N_i$ has a finite irreducible decomposition, giving one for $N$, a contradiction.

It remains to show that an irreducible submodule $Q$ is primary. Pass to $M/Q$. If $xm=0$ with $m\ne0$, the increasing sequence of kernels of multiplication by $x^n$ stabilizes, say at $n$. If $y=x^nm'=am$ belongs to $x^n(M/Q)\cap Am$, then $x^{n+1}m'=xy=0$, so stabilization gives $y=x^nm'=0$. Thus these two submodules intersect in zero. Since zero is irreducible and $Am\ne0$, we have $x^n(M/Q)=0$. Hence every zero-divisor on $M/Q$ is nilpotent.
